\section{Анализ устойчивости}

\subsection{Характеристический полином}
Для системы (4) характеристический полином имеет вид:
\[
\lambda^4 + \left(d + \frac{b}{l}\right)\lambda^3 + \left(c + \frac{g}{l} + \frac{a}{l}\right)\lambda^2 + \frac{g}{l}d\lambda + \frac{g}{l}c = 0. \tag{5}
\]

\subsection{Критерий Рауса-Гурвица}
Для полинома четвертой степени
\[
\lambda^4 + A_3\lambda^3 + A_2\lambda^2 + A_1\lambda + A_0 = 0
\]
с положительными коэффициентами необходимым и достаточным условием устойчивости является:
\[
A_3 A_2 A_1 > A_1^2 + A_3^2 A_0. \tag{6}
\]

\subsection{Условия устойчивости}
Для полинома (5):
\[
A_3 = d + \frac{b}{l} > 0, \quad A_2 = c + \frac{g}{l} + \frac{a}{l} > 0,
\]
\[
A_1 = \frac{g}{l}d > 0, \quad A_0 = \frac{g}{l}c > 0.
\]

Дополнительное условие (6):
\[
\left(d + \frac{b}{l}\right)\left(c + \frac{g}{l} + \frac{a}{l}\right)\left(\frac{g}{l}d\right) > 
\left(\frac{g}{l}d\right)^2 + \left(d + \frac{b}{l}\right)^2 \left(\frac{g}{l}c\right). \tag{7}
\]

\subsection{Численный пример}
Для параметров: \(l = 0.5\) м, \(g = 9.81\) м/с\(^2\), выберем:
\[
a = 15,\; b = 2,\; c = 1,\; d = 1.
\]

Корни характеристического полинома:
\[
\lambda_{1,2} = -0.97 \pm 4.11i,\quad \lambda_{3,4} = -1.03 \pm 1.12i.
\]
Все корни имеют отрицательные действительные части — система устойчива.